\documentclass[11pt]{article}
\usepackage[margin=1in]{geometry}
\usepackage{amsmath,amssymb,mathtools}
\usepackage{fontspec,unicode-math}
\setmainfont{Latin Modern Roman}
\setsansfont{Latin Modern Sans}
\setmathfont{Latin Modern Math}
\usepackage{xurl,hyperref}
\hypersetup{colorlinks=true,urlcolor=blue,linkcolor=blue}
\newcommand{\sref}[2]{\hyperlink{source:#1:#2}{[S#2]}}
\newcommand{\cataloguescope}[1]{\paragraph{Recorded mathematical question.}#1}
\begin{document}
\section{Two-process coupling characterization of weak convergence}\label{statement}

\subsection{Definitions and mathematical statement}

\paragraph{Shared notation.}$\mathbb N=\{1,2,\ldots\}$, and $\mathbb Z,\mathbb Q,\mathbb R,\mathbb C$ denote the integers, rationals, reals and complex numbers. Any other conventions are specified locally.


Let $E$ be a separable metric space with metric $\rho$ and its Borel sigma-field. On $E^{\mathbb N_0}$ use a compatible product metric, for example
\[d(x,y)=\sum_{k=0}^\infty2^{-k-1}\min\{1,\rho(x_k,y_k)\}.\]
Let $X,X'$ be processes in this path space and $\theta_n(x)_k=x_{n+k}$. Suppose $\theta_nX\Rightarrow X'$ as random path-space elements. A two-process coupling constructs only $X$ and $X'$ together, preserving their respective full laws; it does not introduce a separate copy with law $\mathcal L(\theta_nX)$ for each $n$.
\paragraph{Statement recorded in the supplied source.}
Give a characterization of this weak convergence using a two-process coupling. In particular, can such a coupling be chosen so that $d(\theta_nX,\theta_nX')\to0$ almost surely as $n\to\infty$? \sref{statement}{2}
\subsection{Short English statement}
Can weak convergence of shifted processes be characterized by coupling just the original process and its limiting process?
\subsection{Sources}
\begin{description}
\item[\hypertarget{source:statement:1}{S1}] ULAM.AI and the UnsolvedMath Contributors, \emph{UnsolvedMath: A Curated Collection of Open Mathematics Problems}, catalogue entry. CC BY 4.0. \url{https://huggingface.co/datasets/ulamai/UnsolvedMath}.
\item[\hypertarget{source:statement:2}{S2}] Hermann Thorisson, \emph{Some Open Probability Problems}, author preprint; original Universidad de Buenos Aires copy preserved in Internet Archive capture of 27 March 2024, Section 3, weak-convergence paragraph and Problem 3.3, p. 4. \url{https://web.archive.org/web/20240327021031id_/http://cms.dm.uba.ar/depto/public/Some%20Open%20Problems-preprint.pdf}.
\end{description}
\label{end:statement}
\end{document}
